\section{Related Work}
\label{sec:related_work}

\paragraph{Supervision bottlenecks in surgical video understanding.}
Surgical video understanding has often been studied through supervised tasks such as phase recognition, workflow modeling with timestamp labels, and action-triplet recognition~\cite{tecno2020,lessismore2022,rendezvous2022}. These settings have driven substantial progress, but they also highlight the practical bottleneck of surgical AI: obtaining dense expert annotations for phases, steps, tools, actions, and targets is costly, center-specific, and difficult to scale across procedures. This limitation motivates representation learning paradigms that can exploit weak supervision and support more transferable or zero-shot evaluation.

\paragraph{Surgical video-language pretraining.}
Against this background, surgical video-language pretraining has become an important direction for scalable and transferable surgical representation learning. Early work such as SurgVLP established surgical lecture videos as a viable weak supervision source for CLIP-style pretraining~\cite{lecturevideo2023}. Later studies extended this line through hierarchy-aware supervision, procedure-aware augmentation, retrieval-based enhancement, and larger-scale curated corpora, including HecVL, PeskaVLP, OphCLIP, and SurgLaVi~\cite{hecvl2024,peskavlp2024,ophclip2024,surglavi2025}. These works collectively show that weak video-text supervision can improve zero-shot or low-supervision surgical understanding, especially when the textual side becomes larger, cleaner, and more hierarchically structured.

\paragraph{Alignment reliability under weak surgical supervision.}
At the same time, recent surgical VLP papers make clear that richer text does not automatically imply better alignment. SurgVLP notes semantic misalignment between clips and text in surgical lecture videos~\cite{lecturevideo2023}, SurgLaVi highlights noisy pairs caused by temporal misalignment and commentary drift~\cite{surglavi2025}, and PeskaVLP identifies textual information loss and cross-modal procedural alignment as key challenges~\cite{peskavlp2024}. However, most CLIP-style surgical VLP pipelines still optimize standard contrastive learning over selected clip-caption pairs, effectively treating the sampled interval as a single positive unit once the pair has been constructed. This leaves open a more specific question: how to denoise weak temporal positives inside an otherwise valid clip-caption pair.

\paragraph{Partial-positive learning and text-guided evidence selection.}
Outside the surgical domain, weak video-text learning has increasingly been treated as an alignment-quality problem rather than a simple matter of text availability. Multiple-instance learning provides a useful conceptual lens: a weakly labeled bag can be positive even when only a subset of instances truly supports the label~\cite{dietterich1997mil,ilse2018attentionmil}. In video understanding, weakly supervised temporal localization further shows that language can help suppress irrelevant snippets and foreground latent evidence under coarse supervision~\cite{li2023boostingwtal,hu2024mmplateau,li2024fps}. Text-conditioned aggregation methods such as X-Pool similarly show that fixed mean pooling can be replaced by relevance-aware weighting~\cite{gorti2022xpool}. Our window denoising branch is most closely aligned with this line of work: we treat the expanded interval as a bounded partially positive bag and use the paired caption to estimate latent evidence weights during training.

\paragraph{Hierarchy as a granularity prior.}
Hierarchical video-language datasets and models such as SurgLaVi, HiTeA, and HecVL show that temporal granularity should not be ignored during representation learning~\cite{surglavi2025,ye2023hitea,hecvl2024}. In our work, however, hierarchy is not the main supervision target. Instead, we use fine/mid/coarse labels only as a lightweight prior on expected evidence concentration during denoising: fine captions should induce sharper frame weighting, whereas broader captions should permit smoother evidence distributions.

\paragraph{Position of our work.}
Our work is closest to the surgical CLIP-style pretraining line, but differs in what it treats as the unresolved bottleneck. We do not primarily contribute a larger corpus, a new backbone, or a new hierarchical objective. Instead, we focus on how weak temporal supervision is consumed during training once a pair has already been constructed. Specifically, we argue that a globally valid weak pair can still contain temporally noisy positives, and we address that issue with train-time window denoising while preserving standard CLIP-style inference at test time.
